% Table 2 (file name kept from the four-table layout; it prints second because Methods cites it second) -- the
% transitional persons by subtype: what each subtype is, how many persons carry it per source, what truth the subtype
% rests on, where the built split puts them, how many of them the current encoder trained on, the unread COBRA2026
% screen calls, and the best-case Wilson bound their test count allows. Rewritten 2026-09-23 around Figures 1 and 5:
% the claim columns of the earlier table (sensitivity, specificity, exact level, registration ceiling, simulated
% surgeon) are gone; the comparison (formerly Table 3, tab_levelcheck.tex, and Figure 6) is paper 2 since 2026-09-24
% (paper2_levelcheck/). No merged cells. Sources, one line per cell, in tab4_transitional.sources.md; the same rows
% in tab4_transitional.md.
%   results/rare_class/transitional_subtypes.csv  persons per subtype and source (distinct patient_key; 196 rows, 194 persons)
%   results/rare_class/split_plan.md               the built G3 split (make_split.py apply 2026-09-23 20:18, seed 0,
%                                                  leak audit PASS): test persons per subtype, and val+test persons in
%                                                  the v8 or v9 level-bank train split. OWNED BY THE SPLIT WORKFLOW.
%   results/power/inputs/cobra_lstv_screen.grid.json  the COBRA2026 screen calls (SPINEPS+VERIDAH, unread), as Figure 5 draws them
%   figures/figure5.record.json (panel d)          Wilson lower bounds at zero errors; recomputed here (z = 1.959964)
% Updated 2026-09-24 to the rebuilt Figure 5 (figures/figure5.record.json): the COBRA2026 column adds, in brackets,
% the calls their own field supports (57 of 91: L6 11, T13 38, four-lumbar 8), the Wilson column their best case,
% and footnote c said the three real C-ArmLSpine L6 are excluded from the comparison's primary analysis in advance
% (LEVELCHECK_MODEL.md s4), so their 43.9 % bound is no longer printed as an available real-frame bound.
% Revision 5 (2026-09-24, REVISION_LOG3.md section 8): the comparator clauses of notes c, d and f moved to paper 2
% (paper2_levelcheck/MOVED_FROM_PAPER1.md item 3); note f now defines "at risk" itself; note c points to the
% cross-validation row of the naming table.
% Revision 2 (2026-09-24, REVISION_LOG2.md): the COBRA2026 bracket and the "real" Wilson bounds follow the rebuilt
% Figure 5 (figures/figure5.record.json, 08:26; results/transitional_registry/cobra/rows.csv built 06:49): persons a
% two-labeller count (SPINEPS+VERIDAH and a TotalSegmentator rib count, both junctions in the field) calls
% transitional AND that have an exported frame (L6 3, four lumbar 4, T13 0; pooled at risk 7), replacing the 57
% "supported" screen calls; the Any row is the frozen at-risk strata (L6, T13, four-lumbar sacralised: 51 CT test
% persons, Wilson 93.0) with the own-label sacralised L5 as naming unaffected (REVIEW2_STATS N7); footnote e adds the
% labeller's T13 record, 0 of 52 on CTSpinoPelvic1K (results/spineps_ids/IMPACT.md line 230); footnote d drops the
% bank names (REVIEW2_RSNA Fix 26).
% Layout: landscape (rotated 23.2 cm minipage, see tab_corpus.tex); 7 pt on 7.6 pt.
\begin{table}[p]
\centering
\begingroup
\singlespacing
\fontsize{7}{7.6}\selectfont
\newcommand{\fnm}[1]{\textsuperscript{\fontsize{7}{7}\selectfont #1}}
\renewcommand{\arraystretch}{1.0}
\rotatebox[origin=c]{90}{\begin{minipage}{23.2cm}
\captionsetup{font=footnotesize,skip=4pt}
\centering
\caption{Transitional persons by subtype: definition, persons per source, the truth behind the subtype, the test side of the built person-disjoint split, the share of held-out persons the current level encoder trained on, the unread COBRA2026 screen calls, and the best-case Wilson 95\,\% lower bound the test count allows.}
\label{tab:transitional}
% widths sum to 21.90 cm; 14 gaps of 1.9 pt (0.94 cm) -> 22.84 cm
\setlength{\tabcolsep}{1.9pt}
\begin{tabular}{%
>{\raggedright\arraybackslash}p{2.30cm}%
>{\raggedright\arraybackslash}p{4.50cm}%
>{\raggedright\arraybackslash}p{2.70cm}%
>{\raggedright\arraybackslash}p{3.80cm}%
>{\raggedright\arraybackslash}p{2.10cm}%
>{\raggedright\arraybackslash}p{2.30cm}%
>{\raggedright\arraybackslash}p{1.90cm}%
>{\raggedright\arraybackslash}p{2.30cm}}
\toprule
Subtype\fnm{a} & Definition; what a detector can see\fnm{b} & Persons: total (CTSP / VerSe / TS / C-Arm) & Truth behind the subtype & Test persons, built split\fnm{c}: CT + real & Validation and test persons the current encoder trained on\fnm{d} & Unread COBRA2026 screen calls (two-labeller count, with a frame)\fnm{e} & Wilson lower bound with no error at the test count, \%\fnm{f} \\
\midrule
L6 & six rib-free vertebrae above the sacrum; six lumbar bodies where the sixth carries a label & 110 (18 / 86 / 3 / 3) & CTSP two-reader consensus or release label; VerSe published read (38 persons on the audit read alone); TS independent rib count; C-Arm one automatic labeller & 43 + 3 & 49 of 54 & 38 (3) & 91.8 (43 CT); 43.9 (3 real) \\
\addlinespace[2pt]
T13 & thirteen rib-bearing vertebrae & 7 (0 / 6 / 1 / 0) & VerSe mask label 28; TS independent rib count & 3 + 0 & 4 of 4 & 45 (0) & 43.9 (3 CT); none (0 real) \\
\addlinespace[2pt]
Sacralised L5, four lumbar & L5 fused to the sacrum (Castellvi III or IV) and four lumbar labels; the fused body sits in the sacrum label: one box & 12 (9 / 3 / 0 / 0) & CTSP two-reader consensus; VerSe published read & 5 + 0 & 14 of 14, both sacralised rows & 8 (4) & 56.6 (5 CT); 51.0 (4 real) \\
\addlinespace[2pt]
Sacralised L5, five lumbar & L5 fused to the sacrum on a five-lumbar count; its own label (CTSP) or none (VerSe) & 16 (3 / 13 / 0 / 0) & as the row above & 6 + 0 & with the row above & none called & 61.0 (6 CT) \\
\addlinespace[2pt]
Castellvi I/II & enlarged or articulating transverse process; the count is unchanged & 31 (4 / 27 / 0 / 0) & CTSP two-reader consensus; VerSe published read & 12 + 0 & 15 of 15 & not screened & 75.8 (12 CT) \\
\addlinespace[2pt]
Lumbar rib & rib on L1; the count is unchanged & 14 (14 / 0 / 0 / 0) & CTSP release labels & 6 + 0 & 7 of 7 & not screened & 61.0 (6 CT) \\
\addlinespace[2pt]
Unassigned & four lumbar labels, source read normal; no subtype established & 4 (4 / 0 / 0 / 0) & CTSP release labels & 2 + 0 & 3 of 3 & -- & -- \\
\addlinespace[2pt]
L6 and T13 in one spine & none in any truth table & 0 & -- & -- & -- & -- & -- \\
\addlinespace
Any subtype & the rows above & 194 (52 / 135 / 4 / 3) & -- & 77 + 3 & 92 of 97 & 91 (7) & at-risk rows (L6, T13, four lumbar): 93.0 (51 CT); 64.6 (7 real) \\
\bottomrule
\end{tabular}
\par\vspace{1pt}
\begingroup\fontsize{7}{7.6}\selectfont\raggedright
CTSP = CTSpinoPelvic1K; TS = TotalSegmentator; C-Arm = C-ArmLSpine; L5, L6 = fifth, sixth lumbar vertebra; T13 = thirteenth thoracic vertebra.
\fnm{a}Distinct persons (patient key; two VerSe series pairs are one person each); rib convention: L6 = six rib-free vertebrae between the last rib-bearing vertebra and the sacrum, T13 = thirteen rib-bearing vertebrae. Counted 23 September 2026 from the subtype table; eight of twelve TotalSegmentator spines once called L6 from raw mask identifiers read normal on an independent rib count and are not counted.
\fnm{b}Whether a detector sees one body or two is decided by the mask: a fused body with its own label gives two boxes, one inside the sacrum label gives one, one left unlabelled gives none.
\fnm{c}Person-disjoint, subtype-stratified split, 50\,/\,10\,/\,40 train / validation / test within each subtype, every real transitional subject on the test side (the three C-ArmLSpine L6); built 23 September 2026, leak audit passed; Figure~\ref{fig:power}c draws these counts. The three real C-ArmLSpine L6 are scored only in cross-validation, as agreement with one labeller's read (Table~\ref{tab:naming}).
\fnm{d}Persons on the validation or test side of the built split who sit in the training split of the level bank behind the current encoder (the two most recent training sets); a model descending from those banks cannot score them, so a transitional test needs an arm trained without them.
\fnm{e}SPINEPS with VERIDAH on the planning CT deformably registered to the cone-beam CT, read by no person; on CTSP this labeller's L6 call matched the readers in 10 of 24 and its T13 call in 0 of 52. The screen also flagged three cases for re-reading. All calls sit on the test side; none counts as a transitional test person until verified. In brackets, counted 24 September 2026, the persons whom this labeller and an independent TotalSegmentator rib count, both requiring the thoracolumbar and lumbosacral junctions in the field, call transitional and who have an exported frame; no T13 can be counted, T1 lying outside nearly every field; four-lumbar persons are placed in the four-lumbar row.
\fnm{f}Wilson 95\,\% lower bound when every test person is named correctly, the best case at that count (``real'': the bracketed COBRA2026 persons, if each were verified); at risk = L6, T13 and the four-lumbar sacralised L5, the subtypes a count from the sacrum or from C2 misnames; an own-label sacralised L5, Castellvi I/II and a lumbar rib leave the count unchanged and are not in the at-risk bound; the pre-stated targets are 90\,\% for the exact level (35 persons with no error) and 85\,\% for variant sensitivity (22 carriers without a miss).
\par\endgroup
\end{minipage}}
\endgroup
\end{table}
